% ==============================================================================
% MACROS Y ATAJOS MATEMÁTICOS PARA PROBABILIDAD Y ESTADÍSTICA
% Diseñados para escribir fórmulas limpias, legibles y fáciles de modificar
% ==============================================================================

% ------------------------------------------------------------------------------
% 1. PROBABILIDADES Y CONJUNTOS
% ------------------------------------------------------------------------------
% Probabilidad de un evento: \prob{A} -> P(A)
\newcommand{\prob}[1]{\mathbb{P}\left(#1\right)}

% Probabilidad condicional: \probcond{A}{B} -> P(A | B)
\newcommand{\probcond}[2]{\mathbb{P}\left(#1 \;\middle|\; #2\right)}

% Conjunto de resultados o espacio muestral: \espacioMuestral
\newcommand{\espacioMuestral}{\Omega}

% ------------------------------------------------------------------------------
% 2. ESPERANZA, VARIANZA Y MOMENTOS
% ------------------------------------------------------------------------------
% Esperanza matemática: \esp{X} -> E[X]
\newcommand{\esp}[1]{\mathbb{E}\left[#1\right]}

% Esperanza condicional: \espcond{X}{Y} -> E[X | Y]
\newcommand{\espcond}[2]{\mathbb{E}\left[#1 \;\middle|\; #2\right]}

% Varianza: \var{X} -> Var(X)
\newcommand{\var}[1]{\mathrm{Var}\left(#1\right)}

% Desvío estándar: \sd{X} -> \sigma(X)
\newcommand{\sd}[1]{\mathrm{SD}\left(#1\right)}

% Covarianza: \cov{X}{Y} -> Cov(X, Y)
\newcommand{\cov}[2]{\mathrm{Cov}\left(#1, #2\right)}

% Coeficiente de correlación: \corr{X}{Y} -> \rho(X, Y)
\newcommand{\corr}[2]{\rho\left(#1, #2\right)}

% ------------------------------------------------------------------------------
% 3. DISTRIBUCIONES DE PROBABILIDAD USUALES
% ------------------------------------------------------------------------------
% Distribución Binomial: \distBinomial{n}{p} -> Binomial(n, p)
\newcommand{\distBinomial}[2]{\mathrm{Binomial}\left(#1, #2\right)}

% Distribución Poisson: \distPoisson{\lambda} -> Poisson(\lambda)
\newcommand{\distPoisson}[1]{\mathrm{Poisson}\left(#1\right)}

% Distribución Geométrica: \distGeometrica{p} -> Geom(p)
\newcommand{\distGeometrica}[1]{\mathrm{Geom}\left(#1\right)}

% Distribución Hipergeométrica: \distHipergeom{N}{K}{n} -> Hipergeom(N, K, n)
\newcommand{\distHipergeom}[3]{\mathrm{Hipergeom}\left(#1, #2, #3\right)}

% Distribución Uniforme: \distUniforme{a}{b} -> U(a, b)
\newcommand{\distUniforme}[2]{\mathcal{U}\left(#1, #2\right)}

% Distribución Normal: \distNormal{\mu}{\sigma^2} -> N(\mu, \sigma^2)
\newcommand{\distNormal}[2]{\mathcal{N}\left(#1, #2\right)}

% Distribución Exponencial: \distExp{\lambda} -> Exp(\lambda)
\newcommand{\distExp}[1]{\mathrm{Exp}\left(#1\right)}

% ------------------------------------------------------------------------------
% 4. CADENAS DE MARKOV
% ------------------------------------------------------------------------------
% Matriz de transición de un paso: \matrizP
\newcommand{\matrizP}{\mathbf{P}}

% Vector de probabilidades de estado o distribución estacionaria: \vectorPi
\newcommand{\vectorPi}{\boldsymbol{\pi}}

% Estado genérico: \estado{i}
\newcommand{\estado}[1]{\text{Estado } #1}

% ------------------------------------------------------------------------------
% 5. CONFIABILIDAD DE SISTEMAS
% ------------------------------------------------------------------------------
% Función de confiabilidad del sistema: \confiabilidad{t} -> R(t)
\newcommand{\confiabilidad}[1]{R\left(#1\right)}

% Tasa de fallas (hazard rate): \tasaFallas{t} -> \lambda(t)
\newcommand{\tasaFallas}[1]{\lambda\left(#1\right)}

% Tiempo medio hasta la falla: \MTTF
\newcommand{\MTTF}{\mathrm{MTTF}}
